\chapter{Full \(\mathcal{V}^{\mathrm{prim}}\) P3 and the Koszul--Petersson Compatibility on the Discrete Core}
\label{v4-arith:w5-c:full-vprim-P3}

The discrete-core pairing problem requires an injective native
realization in square-integrable automorphic functions and maps
which identify its Koszul pairing with that Hilbert pairing.
The local GNS construction determines the operator normalization,
but does not construct these native maps. Temperedness assumptions
and the continuous-spectrum normalization are separate inputs.
The following transfers retain the maps explicitly.

\section{Native Pairing Maps on a Tempered Cuspidal Core}
\label{v4-arith:w5-c:recap-tempered}

The local construction applies to a native core only after its
pairing maps and adelic assembly are supplied.

\begin{theorem}[conditional tempered cuspidal pairing transfer]
\label{v4-arith:w5-c:thm:tempered-recap}
\ClaimStatusProvedHereConditional. Assume the native pairing
maps of \Cref{v4-arith:kp:hyp:native-pairing-maps} on the
chosen tempered cuspidal core, including its actual ramified
factors and its specified adelic assembly. Then
\begin{equation}
 \langle v,w\rangle^{\mathrm{Pet}}
     =\langle v,cw\rangle^{\mathrm{Koszul}}.
 \label{v4-arith:w5-c:eq:tempered-recap}
\end{equation}
Here $c$ is the assembled native conjugate-linear map. Its
identification with GNS modular conjugation is one of the map
conditions, not a consequence of equality of scalar characters.
\end{theorem}
\begin{proof}
The local calculation is
$B_p(v,c_pw)=B_q(f_pv,J_qf_pw)=h_q(f_pv,f_pw)$.
The assumed native metric maps and adelic assembly turn these
local identities into the displayed global identity. This is
\Cref{v4-arith:kms-gns:thm:ramified-cuspidal-tempered-extension}.
The GNS factor used there is Hilbert--Schmidt, with the stated
density; it has no independent orthogonal phase-coordinate factor.
\end{proof}

\begin{remark}[temperedness and modular domains]
\label{v4-arith:w5-c:rem:baseline}
The modular operator on a represented prime GNS space is
unbounded: $\Delta_qE_{0j}=q^{-j}E_{0j}$ for $0<q<1$.
The GNS Hilbert pairing is nevertheless defined on its full
Hilbert completion, and the modular operator has its explicit
self-adjoint domain. No automorphic temperedness hypothesis is
used in that construction. Temperedness restricts a native
representation domain; it does not give a bounded modular
operator or supply its native pairing maps.
\end{remark}

\section{Non-Tempered Cuspidal Representations: Deligne and Kim--Shahidi}
\label{v4-arith:w5-c:non-tempered}

\subsection{The Non-Tempered Locus}
\label{v4-arith:w5-c:non-tempered:locus}

Cuspidal automorphic representations of \(GL_{2}/\mathbb{Q}\) decompose
as \(\pi=\bigotimes'_{v}\pi_{v}\), with archimedean component \(\pi_{\infty}\)
of one of three types: a holomorphic discrete series \(D_{k}\) of
weight \(k\geq 2\); a Maass cusp form with Laplace eigenvalue
\(\lambda=1/4+r^{2}\) (\(r\in\mathbb{R}\) tempered, \(r\in i\mathbb{R}\)
non-tempered); or a character twist. Finite places carry Satake
parameters \((\alpha_{p},\beta_{p})\) with \(\alpha_{p}\beta_{p}=1\)
(unitarity) and \(\max(|\alpha_{p}|,|\beta_{p}|)=p^{\theta_{p}}\).
Temperedness is the condition \(\theta_{p}=0\) at all \(p\); the
Ramanujan conjecture asserts temperedness for all cuspidal \(\pi\).

\begin{itemize}
\item \textbf{Holomorphic sector.} \(\pi_{\infty}=D_{k}\) for
\(k\geq 2\). Deligne 1974 \emph{Publ.~IHES}~43, La conjecture de Weil
I, gives the Ramanujan bound unconditionally: \(|a_{p}|\leq 2 p^{(k-1)/2}\)
for the Hecke eigenvalue \(a_{p}\) of a holomorphic weight-\(k\) newform,
which is temperedness
\(|\alpha_{p}|=|\beta_{p}|=p^{(k-1)/2}\cdot 1\) after normalisation
(in the analytic normalisation \(\pi\) unitary,
\(|\alpha_{p}|=|\beta_{p}|=1\)). Weight-\(1\) holomorphic forms are
covered by Deligne--Serre 1974 \emph{Ann.~Sci.~ENS.~7}, via the
Artin conjecture for two-dimensional odd Galois representations
attached to weight-\(1\) newforms.
\item \textbf{Maass sector.} \(\pi_{\infty}\) is a Maass cusp form with
complementary-series parameter \(r=it\), \(0<t\leq 1/2\), at \(\infty\);
temperedness at \(\infty\) is \(t=0\). At finite places, non-temperedness
is \(\theta_{p}>0\). The Ramanujan conjecture for Maass forms is open:
the best current bound is Kim--Sarnak 2003 Duke~Math.~J.~119
\(\theta_{p}\leq 7/64\), improving earlier Luo--Rudnick--Sarnak 1995
and Kim--Shahidi 2002 \emph{Ann.~Math.}~155, which used the
symmetric-square and symmetric-fourth-power lifts to bound
\(\theta_{p}\leq 1/9\) via Jacquet--Piatetski-Shapiro--Shalika~1983.
\end{itemize}

\subsection{The Holomorphic Sector and Its Native Comparison}
\label{v4-arith:w5-c:non-tempered:holomorphic}

\begin{theorem}[Koszul--Petersson on a realized holomorphic cuspidal sector]
\label{v4-arith:w5-c:thm:holomorphic-cuspidal}
\ClaimStatusProvedHereConditional. Assume the native pairing maps and assembly of \Cref{v4-arith:kp:hyp:native-pairing-maps} on the stated domain. Let \(\mathcal{V}^{\mathrm{prim}}_{\mathrm{cusp,hol}}\subset\mathcal{V}^{\mathrm{prim}}\)
denote the subspace spanned by holomorphic cuspidal automorphic
representations of \(GL_{2}/\mathbb{Q}\), namely those \(\pi\) with
\(\pi_{\infty}=D_{k}\) for some \(k\geq 1\). For all
\(v,w\in\mathcal{V}^{\mathrm{prim}}_{\mathrm{cusp,hol}}\),
\begin{equation}
\label{v4-arith:w5-c:eq:holomorphic}
\langle v,w\rangle^{\mathrm{Pet}}
\;=\;\langle v,\sigma^{\mathrm{BC}}(w)\rangle^{\mathrm{Koszul}}.
\end{equation}
The pairing identification retains those native realization maps;
the automorphic parameter bounds are separate inputs.
\end{theorem}

\begin{proof}
By Deligne 1974 (weight \(k\geq 2\)) and Deligne--Serre 1974
(weight \(k=1\)), every holomorphic cuspidal representation of
\(GL_{2}/\mathbb{Q}\) is tempered at every finite place:
\(|\alpha_{p}|=|\beta_{p}|=1\) at each \(p\). Temperedness at \(\infty\)
for \(\pi_{\infty}=D_{k}\) is automatic (discrete-series
representations are tempered by Harish-Chandra's definition).
The holomorphic cuspidal sector is therefore a subsector of the
tempered cuspidal sector:
\(\mathcal{V}^{\mathrm{prim}}_{\mathrm{cusp,hol}}\subset\mathcal{V}^{\mathrm{prim}}_{\mathrm{cusp,temp}}\).
\Cref{v4-arith:w5-c:thm:tempered-recap} gives
(\ref{v4-arith:w5-c:eq:holomorphic}) under its native pairing
and assembly hypotheses.
\end{proof}

\begin{remark}[domain of Deligne 1974: holomorphic sector only]
\label{v4-arith:w5-c:rem:deligne}
Deligne 1974 establishes the Ramanujan--Petersson bound for the
Hecke eigenvalues of a holomorphic weight-\(k\) cusp form via the
Riemann hypothesis for \(\ell\)-adic cohomology of a suitable
Kuga--Sato variety. The bound transfers to temperedness of the
associated automorphic representation at every finite prime.
Deligne 1974 applies to holomorphic cuspidal forms; for Maass cusp
forms, no algebraic variety whose Frobenius eigenvalues give the
Hecke eigenvalues is available, and the Ramanujan bound in the Maass
sector is an open problem.
The current state of the art on the Maass sector is Kim--Sarnak
2003 \(\theta\leq 7/64\); this is a sub-convex Ramanujan bound, strictly
weaker than full temperedness.
\end{remark}

\subsection{The Maass Sector: Conditional on Ramanujan}
\label{v4-arith:w5-c:non-tempered:maass}

\begin{theorem}[Koszul--Petersson on the Maass cuspidal sector, conditional]
\label{v4-arith:w5-c:thm:maass-cuspidal}
\ClaimStatusProvedHereConditional. Assume the native pairing maps and assembly of \Cref{v4-arith:kp:hyp:native-pairing-maps} on the stated domain.  \emph{(on the Ramanujan bound for Maass cusp
forms of \(GL_{2}/\mathbb{Q}\))}. Let
\(\mathcal{V}^{\mathrm{prim}}_{\mathrm{cusp,Maass}}\subset\mathcal{V}^{\mathrm{prim}}\)
denote the subspace spanned by Maass cuspidal automorphic
representations of \(GL_{2}/\mathbb{Q}\). Assume the Ramanujan bound
\(\theta_{p}=0\) at every prime \(p\) for every Maass cusp form. For
all \(v,w\in\mathcal{V}^{\mathrm{prim}}_{\mathrm{cusp,Maass}}\),
\begin{equation}
\label{v4-arith:w5-c:eq:maass}
\langle v,w\rangle^{\mathrm{Pet}}
\;=\;\langle v,\sigma^{\mathrm{BC}}(w)\rangle^{\mathrm{Koszul}}.
\end{equation}
\end{theorem}

\begin{proof}
Under Ramanujan for Maass forms, \(\mathcal{V}^{\mathrm{prim}}_{\mathrm{cusp,Maass}}\subset\mathcal{V}^{\mathrm{prim}}_{\mathrm{cusp,temp}}\);
\Cref{v4-arith:w5-c:thm:tempered-recap} then applies.
\end{proof}

\begin{proposition}[what a local parameter bound controls]
\label{v4-arith:w5-c:prop:KS-partial}
If $|\alpha_p|\leq p^\theta$ and $\sigma>\theta$, then
$|(1-\alpha_pp^{-s})^{-1}|\leq(1-p^{\theta-\sigma})^{-1}$
for $\operatorname{Re}s=\sigma$. At $\theta=7/64$ and
$\sigma=1/2$ this denominator is $1-p^{-25/64}$.
A bound on a difference of native pairings further requires their
comparison maps. For bounded maps $f,g$ from a Hilbert space
into a common Hilbert space, one has
\begin{equation}
 |h(fv,fw)-h(gv,gw)|
   \leq\|f-g\|(\|f\|+\|g\|)\|v\|\|w\|.
 \label{v4-arith:w5-c:eq:KS-local}
\end{equation}
No bound on $\|f-g\|$ is supplied by the scalar denominator
estimate alone.
\end{proposition}
\begin{proof}
The reverse triangle inequality gives
$|1-\alpha_pp^{-s}|\geq1-p^{\theta-\sigma}>0$.
For the pairing estimate, subtract and add $h(gv,fw)$,
then apply Cauchy--Schwarz to the two terms. Rescaling a native
bilinear pairing while leaving its Satake parameters fixed shows
why a parameter bound alone cannot control that pairing's
comparison with a given Hermitian norm.
\end{proof}

\begin{remark}[CAP representations and the residual spectrum]
\label{v4-arith:w5-c:rem:CAP-residual}
For \(GL_{2}/\mathbb{Q}\), the residual spectrum \(L^{2}_{\mathrm{res}}\)
consists of one-dimensional automorphic representations lifted from
the determinant character (Moeglin--Waldspurger 1995
\emph{Spectral Decomposition and Eisenstein Series} II.2; specifically
the trivial character gives the one-dimensional residual piece,
measured by the Fock vacuum
\(|0\rangle\in\mathcal{V}^{\mathrm{prim}}\)). CAP representations
(cuspidal-associated to parabolic) do not contribute to
\(GL_{2}/\mathbb{Q}\): they first appear for \(Sp_{4}\) and higher rank
(Piatetski-Shapiro 1983 \emph{Invent.~Math.}~71). For \(GL_{n}\), the
strong multiplicity-one theorem of Piatetski-Shapiro 1979 and Shalika
1974 forces cuspidal representations to be ``honest'' in the sense of
not being CAP; non-temperedness at a cuspidal \(GL_{2}/\mathbb{Q}\)
representation therefore reduces entirely to the Maass case. The
theorem
(\Cref{v4-arith:w5-c:thm:maass-cuspidal}) conditional on Ramanujan
covers the Maass sector, and the residual spectrum (one-dimensional
trivial character) is covered by the vacuum.
\end{remark}

\section{The Eisenstein Density Mismatch and Domain Restriction}
\label{v4-arith:w5-c:eisenstein}

\subsection{The Mismatch}
\label{v4-arith:w5-c:eisenstein:mismatch}

The scalar normalization proposed for the continuum can already be
tested before its native pairing is constructed.

\begin{proposition}[the scalar continuum obstruction]
\label{v4-arith:w5-c:prop:mismatch}
Use \(\Lambda(u)=\pi^{-u/2}\Gamma(u/2)\zeta(u)\), and put
\(M(s)=\Lambda(2s-1)/\Lambda(2s)\).
At a regular point \(s=\tfrac12+it\), where all the displayed
factors are finite and nonzero, define the two proposed coefficients
\begin{align}
\label{v4-arith:w5-c:eq:Pet-density}
\operatorname{lead}^{\mathrm{Pet}}_{\mathrm{cont}}(s)
 &=\bigl(1+|M(s)|^2\bigr)|\Lambda(2s)|^{-2},\\
\label{v4-arith:w5-c:eq:Koszul-density}
\operatorname{lead}^{\mathrm{Koszul}}_{\mathrm{cont}}(s)
 &=\frac{\Lambda(2s)}{\Lambda(2s-1)}|\Lambda(2s)|^{-2}.
\end{align}
Their moduli differ by a factor of two. Thus they cannot be equal.
These are specified scalar candidates; their names do not identify
them with actual native or regularized automorphic measures.
\end{proposition}

\begin{proof}
The functional equation and conjugation give
\(\Lambda(2it)=\Lambda(1-2it)=\overline{\Lambda(1+2it)}\).
Hence \(|M(s)|=1\). The first coefficient has modulus
\(2|\Lambda(2s)|^{-2}\), the second
\(|\Lambda(2s)|^{-2}\). This proves the inequality on the stated
domain. The thermal Euler ratio constructs the second scalar ratio,
as proved in \Cref{v4-arith:kms-gns:lem:Koszul-continued}; it contains
no native representation argument and so cannot supply the missing
pairing or spectral-measure maps.
\end{proof}

\subsection{Domain Restriction}
\label{v4-arith:w5-c:eisenstein:domain}

\begin{definition}[discrete-spectrum core of \(\mathcal{V}^{\mathrm{prim}}\)]
\label{v4-arith:w5-c:def:discrete-core}
The \emph{discrete-spectrum core} is
\begin{equation}
\label{v4-arith:w5-c:eq:discrete-core-def}
\mathcal{V}^{\mathrm{prim}}_{\mathrm{disc}}\;:=\;\mathcal{V}^{\mathrm{prim}}_{\mathrm{cusp}}\;\oplus\;\mathcal{V}^{\mathrm{prim}}_{\mathrm{res}},
\end{equation}
the image of \(\mathcal{V}^{\mathrm{prim}}\) inside
\(L^{2}_{\mathrm{disc}}([GL_{2}])=L^{2}_{\mathrm{cusp}}\oplus L^{2}_{\mathrm{res}}\)
under the placement map
\(\iota\) of \Cref{v4-arith:kms-gns:prop:Vprim-placement}.
Equivalently, the image of \(\mathcal{V}^{\mathrm{prim}}\) obtained by
removing the Eisenstein continuum from its Langlands-decomposition
placement.
\end{definition}

The pairing identity on each summand determines only the diagonal
terms of the global form. To sum these identities, one must also
control the mixed terms.

\begin{definition}[compatible pairing blocks]
\label{v4-arith:w5-c:def:compatible-pairing-blocks}
Let $V=V_1\oplus V_2$ and $V^!=V_1^!\oplus V_2^!$ be complex
vector spaces. Let $h$ be a Hermitian form on $V$, linear in its
first variable, $B:V\times V^!\to\mathbb C$ a bilinear form,
and $c:V\to V^!$ a conjugate-linear map. The decompositions
are \emph{compatible pairing blocks} if, for $i\ne j$,
\[
 h(V_i,V_j)=0,\qquad c(V_i)\subset V_i^!,\qquad
 B(V_i,V_j^!)=0.
\]
These are conditions on the given forms and map. Hermitian
orthogonality alone imposes no condition on the mixed terms of $B$.
\end{definition}

\begin{lemma}[assembly of pairing maps]
\label{v4-arith:w5-c:lem:block-pairing-transfer}
Suppose the decompositions are compatible pairing blocks.
On each block suppose there are complex-linear maps
$f_i:V_i\to\mathscr G_i$ and $g_i:V_i^!\to\mathscr G_i$.
Here $\mathscr G_i$ is a complex Hilbert space with Hermitian form
$h_i$, bilinear form $B_i$, and conjugate-linear map $J_i$ satisfying
$B_i(X,J_iY)=h_i(X,Y)$. Assume
\[
 B(v_i,a_i)=B_i(f_iv_i,g_ia_i),\qquad
 h(v_i,w_i)=h_i(f_iv_i,f_iw_i),\qquad
 g_i(c v_i)=J_i f_iv_i.
\]
Then $B(v,cw)=h(v,w)$ on $V$.
\end{lemma}
\begin{proof}
On $\mathscr G=\mathscr G_1\oplus\mathscr G_2$ define
\[
 B_{\oplus}(X,Y)=\sum_{i=1}^2B_i(X_i,Y_i),\qquad
 h_{\oplus}(X,Y)=\sum_{i=1}^2h_i(X_i,Y_i),\qquad
 J_{\oplus}Y=(J_1Y_1,J_2Y_2).
\]
The maps $f=f_1\oplus f_2$ and $g=g_1\oplus g_2$ have orthogonal
images for distinct blocks. Bilinearity and mixed-term vanishing give
$B(v,a)=B_{\oplus}(fv,ga)$. Hermitian orthogonality gives
$h(v,w)=h_{\oplus}(fv,fw)$. Preservation of the dual blocks gives
$gc=J_{\oplus}f$. Consequently
\[
 B(v,cw)=B_{\oplus}(fv,J_{\oplus}fw)=h_{\oplus}(fv,fw)=h(v,w).
\]
Equivalently, expanding both sides leaves just the two diagonal
terms. These are finite sums. The same proof applies to an
algebraic direct sum of any family of blocks, since each vector
has finite support. Extension to a topological completion requires
that the forms and maps extend continuously on the specified domains.
\end{proof}

\begin{example}[diagonal identities do not determine mixed terms]
\label{v4-arith:w5-c:ex:mixed-pairing-obstruction}
Take $V=V^!=\mathbb C^2$, with coordinate vectors $e_1,e_2$, and put
\[
 h(v,w)=v_1\overline{w_1}+v_2\overline{w_2},\qquad
 c(w)=(\overline{w_1},\overline{w_2}),\qquad
 B(v,a)=v_1a_1+v_1a_2+v_2a_2.
\]
For each coordinate line use $\mathscr G_i=\mathbb C$,
$h_i(z,w)=z\overline w$, $B_i(z,a)=za$, and $J_iw=\overline w$.
The coordinate maps $f_i,g_i$ satisfy all three blockwise identities
of the lemma. The coordinate lines are $h$-orthogonal, and $c$
preserves them. Nevertheless $B(e_1,ce_2)=1$ while $h(e_1,e_2)=0$.
The missing condition is $B(V_1,V_2^!)=0$. This example concerns
abstract pairing data; identifying these data with an arithmetic
pairing would require a separate construction.
\end{example}

\begin{theorem}[pairing identity on a realized discrete core]
\label{v4-arith:w5-c:thm:discrete-core-identity}
\ClaimStatusProvedHereConditional. Suppose the native discrete core
$V=V_{\mathrm{cusp}}\oplus V_{\mathrm{res}}$ and its specified dual
$V^!=V_{\mathrm{cusp}}^!\oplus V_{\mathrm{res}}^!$ are compatible
pairing blocks in the sense of
\Cref{v4-arith:w5-c:def:compatible-pairing-blocks}, for the native
Petersson form $h$, bilinear Koszul form $B$, and map $c$.
Assume the pairing maps of
\Cref{v4-arith:kp:hyp:native-pairing-maps} on each summand and their
specified adelic assembly on that summand. Then
\begin{equation}
 \langle v,w\rangle^{\mathrm{Pet}}
    =\langle v,cw\rangle^{\mathrm{Koszul}}
 \label{v4-arith:w5-c:eq:discrete-core-identity}
\end{equation}
on that core. A Ramanujan-type hypothesis can enlarge a domain
of constructed maps; it does not construct them. The residual
summand requires its own native map, even when a selected
subspace is a single vacuum line.
\end{theorem}
\begin{proof}
The local pairing maps and their assumed adelic assembly give
$B(v_i,cw_i)=h(v_i,w_i)$ for $i=\mathrm{cusp},\mathrm{res}$.
Write $v=v_{\mathrm{cusp}}+v_{\mathrm{res}}$ and similarly for $w$.
Since $c$ preserves the dual blocks and $B$ vanishes between
distinct blocks, $B(v,cw)=\sum_iB(v_i,cw_i)$.
Hermitian orthogonality gives $h(v,w)=\sum_i h(v_i,w_i)$.
The two expressions agree. This is precisely the finite-sum
transfer of \Cref{v4-arith:w5-c:lem:block-pairing-transfer}.
Positivity or norm one of a local GNS vacuum does not identify
it with an automorphic residual vector or control mixed pairings.
\end{proof}

\begin{remark}[Eisenstein placement, domain-restricted]
\label{v4-arith:w5-c:rem:eisenstein-placement}
The identity (\ref{v4-arith:w5-c:eq:discrete-core-identity}) has domain
\(\mathcal{V}^{\mathrm{prim}}_{\mathrm{disc}}\); the domain restriction
places the Eisenstein continuous spectrum
\(L^{2}_{\mathrm{cont}}([GL_{2}])\) outside this domain.
The scalar obstruction of \Cref{v4-arith:w5-c:prop:mismatch}
compares proposed continuum coefficients on the regular critical
line. It does not construct a native Koszul spectral measure.
An extension of the discrete-core identity requires maps to the
continuous-spectrum carrier and a proof that they transport the
actual pairings and regularized measures. The scalar ratio
\(\Lambda(2s)/\Lambda(2s-1)\) supplies none of these maps;
the discrete-spectrum identity retains its stated domain.
\end{remark}

\subsection{Sufficiency for Li's Criterion}
\label{v4-arith:w5-c:eisenstein:Li}

\begin{proposition}[domain restriction for the Li trace]
\label{v4-arith:w5-c:prop:domain-suffices-for-Li}
\ClaimStatusProvedHereConditional \emph{(conditional on the
H-P\(^{\mathrm{Ar}}\) determinant-trace realisation and on the trace map placing
Li's zero sum in the discrete sector)}. Li's coefficients
\(\lambda_{n}=\sum_{\rho}(1-(1-1/\rho)^{n})\)
(\Cref{v4-arith:rh:thm:Li-criterion}) sum over the discrete
non-trivial zeros of \(\zeta\). The non-trivial zeros of \(\zeta\) contributing to the Li sum are
exactly the zeros of \(\xi_{\mathbb{R}}\), all of which are discrete;
the Eisenstein continuous spectrum contributes no term once the
H-P\(^{\mathrm{Ar}}\) trace has been identified with this zero sum.
Eisenstein-series poles in \(\xi\) at \(s=0,1\) are Tate poles
(\Cref{v4-arith:rh:rem:trivial-zero}), cancelled by the factor
\(s(s-1)/2\) in
\(\xi_{\mathbb R}=(1/2)s(s-1)\Lambda_{\zeta}\). The Hilbert--P\'olya
operator \(\Theta_{\xi}\) of \Cref{v4-arith:rh:hyp:HPAr}, if it exists,
has regularised determinant divisor equal to the non-trivial zero
divisor, whose ordinate multiset is discrete.
The regularised trace
\(\mathrm{Tr}_{\mathrm{reg}}(\mathcal L_{n})\) of
\Cref{v4-arith:rh:prop:Li-is-Virasoro} then integrates against
the discrete zero distribution: the Eisenstein continuous spectrum is
not in the domain of the trace.
\end{proposition}

\begin{proof}
Li's sum is over non-trivial zeros of \(\zeta\), equivalently zeros of
\(\xi_{\mathbb R}\). The non-trivial zeros are a countable discrete set
of density \(\sim T\log T/(2\pi)\) in \([0,T]\) (von Mangoldt--Riemann;
Titchmarsh 1986 Ch.~II). No Eisenstein-series pole contributes to
this set: Eisenstein poles at \(s=0,1\) in \(\Lambda_{\zeta}\) are
exactly the Tate poles eliminated by the factor \(s(s-1)/2\) in the
Li normalisation
(\Cref{v4-arith:rh:rem:trivial-zero}). The zero divisor of
\(\xi_{\mathbb R}\) is therefore discrete. Under the
H-P\(^{\mathrm{Ar}}\) hypothesis it identifies with the determinant
divisor and Weil--Connes trace of \(\Theta_{\xi}\)
(\Cref{v4-arith:rh:hyp:HPAr}). The Li-transform trace
\(\mathrm{Tr}_{\mathrm{reg}}(\mathcal L_{n})\) equals the sum over
this discrete distribution only after that determinant-trace
identification.
\end{proof}

\begin{corollary}[conditional Li path on the discrete core]
\label{v4-arith:w5-c:cor:Li-path-closes}
\ClaimStatusProvedHereConditional. Assume all hypotheses of
\Cref{v4-arith:w5-c:thm:discrete-core-identity}, including compatible
pairing blocks and the native maps with adelic assembly.
\emph{(conditional on Ramanujan for Maass
forms at the cuspidal Maass boundary, on the H-P\(^{\mathrm{Ar}}\)
determinant-trace realisation, and on the Tamagawa-regularised assembly;
conditional on the native maps on the holomorphic cuspidal and residual sectors)}. The sufficient RH path of \Cref{v4-arith:rh:thm:sufficient-path} receives
a conditional P3 input from \(\mathcal{V}^{\mathrm{prim}}_{\mathrm{disc}}\)
(\Cref{v4-arith:w5-c:thm:discrete-core-identity}); the Li criterion
draws on discrete zero sums (\Cref{v4-arith:w5-c:prop:domain-suffices-for-Li})
from the same domain, and the Eisenstein continuum is outside the
domain of both inputs.
\end{corollary}

\begin{proof}
P3 on \(\mathcal{V}^{\mathrm{prim}}_{\mathrm{disc}}\) is the conditional
identity of \Cref{v4-arith:w5-c:thm:discrete-core-identity}. The Li path of
\Cref{v4-arith:rh:thm:VB-implies-F-RH} uses \(\lambda_{n}\) which
are discrete-zero sums after the H-P\(^{\mathrm{Ar}}\) spectral
identification (\Cref{v4-arith:w5-c:prop:domain-suffices-for-Li}).
The Eisenstein continuum is outside the resulting trace domain.
\end{proof}

\section{Langlands--Shahidi Normalisation}
\label{v4-arith:w5-c:LS}

The Koszul--Petersson identity can be extended to the full
\(\mathcal{V}^{\mathrm{prim}}\), including the Eisenstein continuum,
only after the Koszul pairing on principal-series inductions is
identified with the positive Langlands--Shahidi Plancherel measure.
The scalar \(N^{\mathrm{LS}}\) of \Cref{v4-arith:w5-c:def:LS-factor}
is not by itself such an identification. The discrete-core reduction
of Section~\ref{v4-arith:w5-c:eisenstein} is the construction used for the
Li trace; the present section records the conditional continuum
extension.

\subsection{The Langlands--Shahidi Normalisation Factor}
\label{v4-arith:w5-c:LS:factor}

\begin{definition}[Langlands--Shahidi correction factor]
\label{v4-arith:w5-c:def:LS-factor}
The \emph{Langlands--Shahidi normalisation factor} on the
principal-series induction \(\mathrm{Ind}_{B}^{G}(\chi_{s})\) is
\begin{equation}
\label{v4-arith:w5-c:eq:LS-factor}
N^{\mathrm{LS}}(s)\;:=\;\frac{\Lambda(2s-1)}{\Lambda(2s)}\qquad(\mathrm{Re}(s)=1/2).
\end{equation}
At regular critical-line points, \(|N^{\mathrm{LS}}(s)|=1\). The scalar coming from
the Maass--Selberg truncation is part of the Plancherel convention, not
part of the Langlands--Shahidi normalising operator.
\end{definition}

\begin{proposition}[Koszul normalisation absorption]
\label{v4-arith:w5-c:prop:LS-absorption}
\ClaimStatusConditional \emph{(on the identification of the bar--cobar
Koszul spectral measure on the Eisenstein continuum with the Shahidi
2010 normalised Eisenstein measure; see \Cref{v4-arith:w5-c:rem:LS-gap})}. Under the substitution
\(\langle\,\cdot\,,\,\cdot\,\rangle^{\mathrm{Koszul}}\to\langle\,\cdot\,,\,\cdot\,\rangle^{\mathrm{Koszul},\mathrm{LS}}=N^{\mathrm{LS}}(\cdot)\cdot\langle\,\cdot\,,\,\cdot\,\rangle^{\mathrm{Koszul}}\)
on the Eisenstein continuum and the independent identification of the
resulting quadratic form with the Shahidi \(L^{2}\)-normalised measure,
the corrected Koszul normalisation is transported to the
Maass--Selberg Plancherel normalisation:
\begin{equation}
\label{v4-arith:w5-c:eq:LS-match}
\mathsf{Sh}_{\mathrm{LS}}\!\left(\operatorname{lead}^{\mathrm{Koszul}}_{\mathrm{cont}}(s)\right)
\;=\;\operatorname{lead}^{\mathrm{Pet}}_{\mathrm{cont}}(s).
\end{equation}
Here \(\mathsf{Sh}_{\mathrm{LS}}\) denotes the Shahidi Plancherel
transport. No raw equality of scalar densities is asserted before this
transport is supplied.
\end{proposition}

\begin{proof}
On \(\mathrm{Re}(s)=1/2\), \(|M(s)|=1\) (Moeglin--Waldspurger \S IV.3.10)
gives \(|N^{\mathrm{LS}}(s)|=1\). Hence the Langlands--Shahidi factor
changes the scattering phase but does not by itself turn the Koszul
coefficient into a positive Plancherel density. The equality
(\ref{v4-arith:w5-c:eq:LS-match}) is exactly the additional
identification of the bar--cobar Eisenstein measure with the Shahidi
2010 \(L^{2}\)-normalised Eisenstein measure (Shahidi 2010
Thm~3.5.1). The scalar \(1+|M(s)|^{2}\) belongs to the
Maass--Selberg truncation convention and is not inserted into
\(N^{\mathrm{LS}}\).
\end{proof}

\begin{remark}[Koszul--Shahidi measure identification]
\label{v4-arith:w5-c:rem:LS-gap}
The ratio \(\Lambda(2s)/\Lambda(2s-1)\) is a non-constant unit-modulus function
on \(\mathrm{Re}(s)=1/2\); it equals \(1/M(s)\) with \(|M(s)|=1\).
The equality of Plancherel normalisations holds only when the Koszul
Eisenstein measure is identified with the Shahidi 2010
\(L^{2}\)-normalised measure. Whether the bar--cobar Koszul definition
(Beilinson--Drinfeld 2004 \S3.5) assigns exactly this normalisation to
the Eisenstein sector is an open problem. This gap does not affect
\Cref{v4-arith:w5-c:thm:main}: the main theorem uses the domain
restriction of Section~\ref{v4-arith:w5-c:eisenstein}, not the
Langlands--Shahidi extension.
\end{remark}

\begin{theorem}[full-\(\mathcal{V}^{\mathrm{prim}}\) Koszul--Petersson compatibility via Langlands--Shahidi]
\label{v4-arith:w5-c:thm:full-vprim-LS}
\ClaimStatusProvedHereConditional. Assume the discrete-core hypotheses
of \Cref{v4-arith:w5-c:thm:discrete-core-identity}. On the stated
common core, assume also compatible pairing blocks between the
discrete and Eisenstein sectors for the corrected bilinear form
and the given map $c$, as in
\Cref{v4-arith:w5-c:def:compatible-pairing-blocks}.
Retain the native pairing maps and their specified adelic assembly.
\emph{(on the Ramanujan bound for Maass cusp
forms of \(GL_{2}/\mathbb{Q}\); on the identification of the bar--cobar
Koszul spectral measure on the Eisenstein sector with the Shahidi 2010
\(L^{2}\)-normalised Eisenstein measure
(\Cref{v4-arith:w5-c:rem:LS-gap}); and on the Langlands--Shahidi
correction \(N^{\mathrm{LS}}\) as a canonical ingredient of the arithmetic
Koszul pairing on principal-series inductions)}. Assume Ramanujan for
Maass forms. For all $v,w$ in this common core of
$\mathcal V^{\mathrm{prim}}$,
\begin{equation}
\label{v4-arith:w5-c:eq:full-vprim-LS}
\langle v,w\rangle^{\mathrm{Pet}}
\;=\;\langle v,\sigma^{\mathrm{BC}}(w)\rangle^{\mathrm{Koszul},\mathrm{LS}},
\end{equation}
with \(\langle\,\cdot\,,\,\cdot\,\rangle^{\mathrm{Koszul},\mathrm{LS}}\)
the Langlands--Shahidi-corrected Koszul pairing of
\Cref{v4-arith:w5-c:prop:LS-absorption}. The identification extends
the discrete-core identity
(\ref{v4-arith:w5-c:eq:discrete-core-identity}) to the Eisenstein
continuum.
\end{theorem}

\begin{proof}
On the discrete core, \(N^{\mathrm{LS}}|_{\mathrm{disc}}=1\)
(the correction factor is supported on the Eisenstein continuum),
so \Cref{v4-arith:w5-c:thm:discrete-core-identity} applies without
modification. On the Eisenstein continuum,
\Cref{v4-arith:w5-c:prop:LS-absorption} gives the conditional
Plancherel-normalisation match.
The specified adelic assembly gives these identities on each of
the two sectors. The assumed preservation of their dual blocks
and vanishing of the mixed bilinear terms permit their sum,
by \Cref{v4-arith:w5-c:lem:block-pairing-transfer}. Thus the
identity holds on the stated common core. Its extension to a
chosen completion still requires continuity of both pairings
and of $c$ on that completion.
\end{proof}

\begin{remark}[domain of the Langlands--Shahidi correction]
\label{v4-arith:w5-c:rem:LS-cost}
The full \(\mathcal{V}^{\mathrm{prim}}\) identity
(\ref{v4-arith:w5-c:eq:full-vprim-LS}) extends the domain of the
Koszul pairing to the Eisenstein sector only under the measure
identification of \Cref{v4-arith:w5-c:rem:LS-gap}; it is not obtained
by the bare insertion of
\(N^{\mathrm{LS}}\). The factor \(N^{\mathrm{LS}}\) is the
Langlands--Shahidi normalised intertwining operator
(Shahidi 2010 Ch.~2--3; Moeglin--Waldspurger \S IV.3), the natural
normalisation for the Eisenstein functional equation, and its insertion
is a definitional addition to the arithmetic Koszul pairing of
Chapter~\ref{v4-arith:arakelov-bar}, which uses the Bost--Connes
involution alone. The domain restriction of
Section~\ref{v4-arith:w5-c:eisenstein} requires no such addition and
is sufficient only for the discrete trace normalisation of the Li path
under the hypotheses of
\Cref{v4-arith:w5-c:prop:domain-suffices-for-Li}.
\end{remark}

\section{Full \(\mathcal{V}^{\mathrm{prim}}\) P3 on the Discrete Core}
\label{v4-arith:w5-c:P3-statement}

\begin{theorem}[positive pairing on a realized discrete core]
\label{v4-arith:w5-c:thm:P3-discrete-core}
\ClaimStatusProvedHereConditional. In addition to the hypotheses
of \Cref{v4-arith:w5-c:thm:discrete-core-identity}, suppose the
native core has an injective square-integrable realization $I$
which computes its specified Petersson form. Retain, in particular,
the preservation and mixed-term vanishing of its pairing blocks.
Then
\begin{equation}
 h(v,w)=\int (Iv)(g)\overline{(Iw)(g)}\,dg
 \label{v4-arith:w5-c:eq:P3-form}
\end{equation}
is positive definite and satisfies
\begin{equation}
 h(v,w)=B(v,cw).
 \label{v4-arith:w5-c:eq:P3-compatibility}
\end{equation}
\end{theorem}
\begin{proof}
The integral of $|Iv|^2$ is nonnegative. If it is zero, the
$L^2$ vector $Iv$ vanishes; injectivity gives $v=0$.
The equality with the native bilinear pairing is the transfer
of the preceding theorem. This argument does not produce the
map $I$ from local GNS positivity.
\end{proof}

\begin{corollary}[the native positivity input]
\label{v4-arith:w5-c:cor:P3-input}
\ClaimStatusProvedHereConditional. A positivity argument may use
the form of \Cref{v4-arith:w5-c:thm:P3-discrete-core} only
with its native pairing, square-integrable realization, compatible
pairing blocks, and assembly hypotheses. A Li or determinant-trace argument also
retains its separate trace-comparison hypothesis.
\end{corollary}
\begin{proof}
These are exactly the maps used in the proof of positivity and
pairing transfer. The theorem asserts no identification between
an automorphic squared norm and the zero-divisor trace.
\end{proof}

\section{Koszul--Petersson Compatibility on the Domain-Restricted Discrete Core: Main Theorem}
\label{v4-arith:w5-c:main-theorem}

\begin{theorem}[conditional discrete-core Koszul--Petersson comparison]
\label{v4-arith:w5-c:thm:main}
\ClaimStatusProvedHereConditional. On a discrete core satisfying
all hypotheses of \Cref{v4-arith:w5-c:thm:discrete-core-identity},
including compatible pairing blocks, native realization maps,
and their adelic assembly,
\begin{equation}
 \langle v,w\rangle^{\mathrm{Pet}}
       =\langle v,cw\rangle^{\mathrm{Koszul}}.
 \label{v4-arith:w5-c:eq:main}
\end{equation}
The comparison concerns the given core and maps. It does not
identify an unconstructed global Hochschild trace with the
represented prime GNS carrier.
\end{theorem}
\begin{proof}
The equality is the direct-sum and adelic assembly of the proved
local pairing-map identities, as established in
\Cref{v4-arith:w5-c:thm:discrete-core-identity}.
\end{proof}

\begin{remark}[variance and native maps]
\label{v4-arith:w5-c:rem:source-disjointness-comparison}
Petersson is Hermitian and the Koszul form is bilinear. The
conjugate-linear native map $c$ gives the second expression its
Hermitian variance. On the represented model this is
$B_q(X,J_qY)=h_q(X,Y)$, with two distinct maps for a native
factor and its dual. Inserting $J_q$ into a Hermitian pairing
instead would be bilinear and would fail the scalar test.
\end{remark}

\section{Native and Analytic Comparison Conditions}
\label{v4-arith:w5-c:remaining-conditional}

The first native requirement is the construction of pairing
maps, with a compatible adelic domain. Automorphic bounds then
control which representation factors belong to a domain where
those maps have been constructed. They do not establish the
maps, their norm discrepancy or their interaction with phases.
The local modular operator is unbounded independently of those
bounds, while its Hilbert pairing is positive by construction.

Continuous-spectrum pairings require their actual measure and
regularization. Excluding that spectrum defines a smaller
domain but does not construct a native embedding into the
remaining discrete spectrum. A proposed scalar normalization
extends the pairing identity only after its compatibility with
the native maps and spectral measure has been proved.

Finally, a Hilbert--P\'olya or Li-positivity argument needs a
map from the relevant arithmetic test space to the realized
core and a trace or norm identity on that domain. The local
GNS normalization and the radial shell pairing model do not
supply this last map. These missing comparisons are independent
of the explicit positivity of a squared $L^2$ norm.

\section{Bibliography}
\label{v4-arith:w5-c:bibliography}

The primary sources enter in the following order.

\begin{description}
\item[Deligne, P.]
\emph{La conjecture de Weil~I}, Publ.~Math.~IHES~43 (1974), 273--307.
Provides the Ramanujan bound for holomorphic cuspidal newforms of
weight \(k\geq 2\) via the Riemann hypothesis for varieties over finite
fields.
\item[Deligne, P.]
\emph{La conjecture de Weil~II}, Publ.~Math.~IHES~52 (1980), 137--252.
The mixed-sheaf generalisation; the holomorphic Ramanujan bound already
follows from Weil~I above.
\item[Deligne, P. and Serre, J.-P.]
\emph{Formes modulaires de poids~\(1\)}, Ann.~Sci.~ENS.~7 (1974), 507--530.
Weight-\(1\) Ramanujan via Artin for two-dimensional odd Galois
representations.
\item[Kim, H. and Shahidi, F.]
\emph{Functorial products for \(GL_{2}\times GL_{3}\) and the symmetric
cube for \(GL_{2}\)}, Ann.~Math.~155 (2002), 837--893. Symmetric-cube
and symmetric-fourth-power functorial lifts establishing the
\(\theta\leq 1/9\) sub-convex Ramanujan bound for Maass cusp forms of
\(GL_{2}/\mathbb{Q}\).
\item[Kim, H. and Sarnak, P.]
\emph{Refined estimates towards the Ramanujan and Selberg conjectures},
appendix to Kim, \emph{Functoriality for the exterior square of
\(GL_{4}\) and the symmetric fourth of \(GL_{2}\)}, J.~Amer.~Math.~Soc.~16
(2003), 139--183. Improvement to \(\theta\leq 7/64\).
\item[Sarnak, P.]
\emph{Some applications of modular forms}, Cambridge Tracts~99 (1990),
and \emph{Selberg's eigenvalue conjecture}, Notices AMS~42 (1995),
1272--1277. The Selberg trace-formula residual spectrum computation
fixing the Maass--Selberg normalisation on the Eisenstein continuum.
\item[Langlands, R. P.]
\emph{On the Functional Equations Satisfied by Eisenstein Series},
Lecture Notes in Math.~544, Springer 1976 (manuscript 1966).
Spectral decomposition, Maass--Selberg relations, intertwining
operators.
\item[Moeglin, C. and Waldspurger, J.-L.]
\emph{Spectral Decomposition and Eisenstein Series: A Paraphrase of
the Scriptures}, Cambridge Tracts 113, Cambridge UP 1995. Modern
\(GL_{n}\) exposition of Langlands 1966; normalised intertwining
operator.
\item[Shahidi, F.]
\emph{Eisenstein Series and Automorphic \(L\)-Functions}, AMS
Colloquium~58, 2010. The Langlands--Shahidi method and
normalisation.
\item[Gelbart, S.]
\emph{Automorphic Forms on Adele Groups}, Ann.~Math.~Studies~83,
Princeton 1975. \(GL_{2}\) spectral decomposition.
\item[Iwaniec, H.]
\emph{Spectral Methods of Automorphic Forms} (second edition), AMS
Graduate Studies~53, 2002. Maass cusp forms and Petersson
positive-definiteness.
\item[Takesaki, M.]
\emph{Tomita's Theory of Modular Hilbert Algebras and its
Applications}, Lecture Notes in Math.~128, Springer 1970; and
\emph{Theory of Operator Algebras~I}, Springer 1979. GNS
construction, modular conjugation, polar decomposition.
\item[Bost, J.-B. and Connes, A.]
\emph{Hecke algebras, type III factors and phase transitions with
spontaneous symmetry breaking in number theory}, Selecta Math.~1
(1995), 411--457. The Bost--Connes \(C^{*}\)-algebra and KMS\({}_{1}\)
state.
\item[Connes, A. and Marcolli, M.]
\emph{Noncommutative Geometry, Quantum Fields and Motives}, AMS
Colloquium~55, 2008. Modern synthesis.
\item[Casselman, W. and Shalika, J.]
\emph{The unramified principal series of \(p\)-adic groups. II. The
Whittaker function}, Compositio Math.~41 (1980), 207--231. Whittaker
uniqueness at supercuspidals.
\item[Bushnell, C. and Henniart, G.]
\emph{The Local Langlands Conjecture for \(GL(2)\)}, Springer 2006.
Tame and wild parametrisation of supercuspidals.
\item[Godement, R. and Jacquet, H.]
\emph{Zeta Functions of Simple Algebras}, Lecture Notes in Math.~260,
Springer 1972. Local functional equation and \(\varepsilon\)-factors.
\item[Jacquet, H. and Langlands, R. P.]
\emph{Automorphic Forms on \(GL(2)\)}, Lecture Notes in Math.~114,
Springer 1970. Whittaker model, ramified principal series.
\item[Jacquet, H., Piatetski-Shapiro, I., and Shalika, J.]
\emph{Rankin--Selberg convolutions}, Amer.~J.~Math.~105 (1983),
367--464. Local factors and functional equation.
\item[Tate, J.]
\emph{Fourier analysis in number fields and Hecke's zeta functions},
Ph.D.~thesis, Princeton 1950; reprinted Cassels--Fr\"ohlich 1967
Ch.~XV. Restricted tensor product.
\item[Weil, A.]
\emph{Basic Number Theory}, Springer 1967. Tamagawa normalisation.
\item[Selberg, A.]
\emph{Harmonic analysis and discontinuous groups in weakly symmetric
Riemannian spaces with applications to Dirichlet series}, J.~Indian
Math.~Soc.~20 (1956), 47--87. Petersson positive-definiteness.
\item[Arthur, J.]
\emph{Eisenstein series and the trace formula}, Proc.~Symp.~Pure
Math.~33 (1979); and \emph{A trace formula for reductive groups.~I},
Duke Math.~J.~45 (1978), 911--952. Adelic truncation; continuous
spectrum distributional pairing.
\item[Piatetski-Shapiro, I.]
\emph{On the Saito--Kurokawa lifting}, Invent.~Math.~71 (1983),
309--338. CAP representations: first appear at \(Sp_{4}\), absent for
\(GL_{2}\).
\item[Gelfand, I. M., Graev, M. I., and Piatetski-Shapiro, I. I.]
\emph{Representation Theory and Automorphic Functions}, Saunders
1969. Archimedean classification of cuspidal representations.
\end{description}
